% Appendix C.1 draft. Numbers are from `python -m analysis.paper_visit_order` (its printed summary);
% the figure and table it writes are \input below. Single-layer gemma-2-2b only.
\subsection{Does the attribution order matter?}
\label{app:visit-order}

CLCD-elimination visits the latent pool from the weakest to the strongest attribution score (Section~\ref{sec:clcd-search}). To test whether this order matters, we rerun it on the single-layer gemma-2-2b organisms with the pool visited in a uniformly random order, three permutations per organism, which uses no attribution at all. Besides the verified circuit, we report the latents elimination keeps (the \emph{kept set}), which isolates the effect of the order, since the verified circuit also depends on the prefix sweep that adds removed latents back and rounds the size up to the grid.

For \topklora, the random order enlarges every circuit (Figure~\ref{fig:visit-order}, Table~\ref{tab:visit-order}): the kept set triples, from a median of 18 to 54 latents, and the verified circuit grows from 22\% to 61\% of the adapter. Dense circuits change far less: their kept set grows from 165 to 211 latents, and the verified circuit, which must span about 90\% of the adapter under either order, keeps its size in 12 of 15 runs. Every random-order circuit still passes verification, with at most one held-out backdoor response in 4{,}000, so attribution makes \topklora\ circuits small rather than valid. Circuits are also not unique: \topklora\ kept sets from different orders overlap by only 7--23\% (Jaccard index), several times chance, around a shared core of 3--9 latents, so a verified circuit is a minimal circuit rather than the unique one.

\input{figures/figC1_visit_order.tex}
\input{tables/tabC1_visit_order.tex}
